% JAC1964-20261121; /root; GPT-6 Astra Pro.
% A standard spreading argument in the exact axis domain of H9-T7.
\subsection*{H9-T8b. The remaining fixed-axis sign can be tested over finite fields}

\begin{proposition}[H9-T8b]
Fix a prime $p$ and an integer $m\ge2$. For a smooth
codimension-two ample axis $\kappa:W\hookrightarrow A$ of a
genuine Lefschetz pencil on an abelian $2m$-fold in
characteristic $p$, put
\[
 V_W=\ker\bigl(\kappa_*:\operatorname{Alg}^{m-1}(W)
                    \longrightarrow H^{2m+2}(A)(m+1)\bigr),
 \qquad q_W(w)=(-1)^{m-1}\deg_W(w^2).
                                                        \tag{H9.8b.1}
\]
Here algebraic classes and the kernel are taken in one fixed
invertible-prime $\ell$-adic realization, as in H9-T7.
The assertion $q_W(w)>0$ for every nonzero $w\in V_W$ and
every algebraically closed ground field of characteristic $p$
is equivalent to the same assertion for axes over
$\overline{\mathbf F}_p$. In the latter assertion each test
axis and each test cycle may be defined over a finite field.
The analogous equivalence holds with the weak inequality
$q_W(w)\ge0$.

Consequently, under the accepted 032 input, the remaining fixed-axis sign and exact
radical assertion in H9-T7 reduces to this finite-field axis
test. This does not prove the finite-field test or replace its
axis by an abelian variety.
\end{proposition}

\begin{proof}
Only the implication from finite fields to general fields needs
proof. Work over an algebraically closed field $k$ of
characteristic $p$, and fix a rational algebraic class
$0\ne w\in V_W$. Choose one rational cycle $z$ representing it.
The abelian variety, ample line, two pencil sections, and the
finitely many components and coefficients of $z$ descend to
a finitely generated field $k_0/\mathbf F_p$.
Spread them over a connected smooth integral scheme $S$ of
finite type over $\mathbf F_p$ with function field $k_0$.
After replacing $S$ by a nonempty open, the abelian scheme
$\mathscr A/S$, its smooth axis $\mathscr W/S$, and the pencil
have the required properties on every geometric fibre.
One may retain the genuine Lefschetz-pencil condition because
the transversality and ordinary critical-point conditions are
open for this fixed finite-presentation family. The closures
of the finitely many cycle components can also be made flat
over $S$, with the required codimension.

The cycle class of this closure gives a section
\[
 c_z\in H^0\bigl(S,R^{2m-2}f_*\mathbf Q_\ell(m-1)\bigr),
          \qquad f:\mathscr W\longrightarrow S.
                                                        \tag{H9.8b.2}
\]
For precision, the total space $\mathscr W$ is smooth after
the preceding choice of $S$; take the usual cycle class there
and its edge map to the displayed higher direct image.
Compatibility of cycle classes with fibre pullback identifies
its geometric stalk with the class of the fibre cycle.
Smooth proper base change makes this higher direct image
lisse. The relative Gysin map to
$R^{2m+2}(\mathscr A/S)_*\mathbf Q_\ell(m+1)$ sends $c_z$
to a section whose geometric generic stalk is zero. A section
of a lisse sheaf on a connected base is determined by any
geometric stalk. The Gysin image is therefore zero at every
geometric point. For the same reason, the nonzero generic
stalk of $c_z$ remains nonzero at every geometric point.

Choose a closed point $s\in S$. Its residue field is finite.
After algebraic closure, the fibre cycle has class
\[
                    0\ne w_s\in V_{W_s}.
                                                        \tag{H9.8b.3}
\]
Cup product and the trace map commute with smooth proper
specialization. Hence the rational intersection number is
unchanged:
\[
                  (-1)^{m-1}\deg_W(z^2)
                    =(-1)^{m-1}\deg_{W_s}(z_s^2).
                                                        \tag{H9.8b.4}
\]
The hypothesized finite-field inequality on the right proves
the required inequality on the left, including strictness
when requested. No algebraicity theorem for all classes on
the special fibre is used. Extra algebraic classes may appear
there without affecting this argument.

This argument applies to each rational test class individually;
it makes no assumption that all algebraic classes have a
common model or that a chosen rational basis is independent
after extending scalars to $\mathbf Q_\ell$. Conversely a
finite-field axis is already in the general class under
consideration. Finally every finite list of cycles over
$\overline{\mathbf F}_p$ is defined over a finite subfield.
This proves the equivalences. Insert the strict axis
inequality into the orthogonal decomposition H9-T7 and use the accepted 032 input
on the ambient abelian summand to obtain the stated height
and radical reduction.
\end{proof}

\paragraph{Exact source and programme scope.}
This is the usual spreading and specialization argument,
applied to the particular Gysin kernel required by cote44,
scan4, Th.5, with the primitive sign of scan8 and cote16,
no.7, scans34--40. Milne, \emph{Polarizations and
Grothendieck's standard conjectures}, Ann.\ of Math.\ 155
(2002), Theorem~3.3, proof on printed pp.607--608, uses the
same specialization of primitive algebraic classes and their
pairings for abelian varieties. The present argument spells
out why the extra relative Gysin-kernel condition persists.
Smooth proper base change is directly stated in Milne,
\emph{Lectures on Etale Cohomology}, version~2.21,
Theorem~20.2, p.127, with stalk specialization in
Proposition~20.3, p.128. The functorial ring cycle map is
Section~23, Theorem~23.4 and the preceding properties(a),(b),
p.142. This is an application of the SGA base-change and
cycle-class framework, not a new theorem attributed to an
unverified SGA/EGA locator.
